\documentclass[11pt]{article}

% Shared notes settings for Elements of AIML lectures (UPES).
% Keep this file free of \documentclass to allow reuse via \input.

\usepackage[utf8]{inputenc}
\usepackage[T1]{fontenc}
\usepackage[a4paper,margin=1in]{geometry}
\usepackage{hyperref}
\usepackage{xurl}
\usepackage{booktabs}
\usepackage{amsmath}
\usepackage{amssymb}
\usepackage{listings}
\usepackage{xcolor}
\usepackage{enumitem}
\usepackage{parskip}

% Code listings (general-purpose).
\lstset{
  basicstyle=\ttfamily\small,
  keywordstyle=\color{blue},
  commentstyle=\color{gray},
  stringstyle=\color{teal},
  showstringspaces=false,
  columns=fullflexible,
  frame=single,
  framerule=0.2pt,
  breaklines=true
}

\setlist[itemize]{topsep=0.3em,itemsep=0.2em}
\setlist[enumerate]{topsep=0.3em,itemsep=0.2em}

% Simple per-section source line.
\newcommand{\notesources}[1]{\par\noindent{\small\color{gray}\textit{Sources:} \sloppy #1\par}}

% Unit-wise source strings for this subject (used by slides/notes).
% Path is relative to the lecture `latex/` folder (the compilation CWD).
\IfFileExists{../../../_shared/aiml-sources.tex}{%
  % Source strings for Elements of AIML (used by slides + notes).

\newcommand{\AIMLRepoURL}{https://github.com/tali7c/Elements-of-AIML}

\newcommand{\AIMLLabSources}{%
Provost and Fawcett, \textit{Data Science for Business}.;%
McKinney, \textit{Python for Data Analysis}.;%
WEKA Documentation (Waikato Environment for Knowledge Analysis), accessed 2026-02-28.;%
Matplotlib and Seaborn documentation, accessed 2026-02-28.%
}

\newcommand{\AIMLUnitISources}{%
Rich and Knight, \textit{Artificial Intelligence}, McGraw Hill, 2017.;%
Russell and Norvig, \textit{Artificial Intelligence: A Modern Approach} (4e), Ch. 1--2 (Introduction, Intelligent Agents).;%
Poole and Mackworth, \textit{Artificial Intelligence: Foundations of Computational Agents} (2e), selected sections.%
}

\newcommand{\AIMLUnitIISources}{%
Russell and Norvig, \textit{Artificial Intelligence: A Modern Approach} (4e), Ch. 7--9 (Logical Agents, First-Order Logic, Inference).;%
Huth and Ryan, \textit{Logic in Computer Science} (2e), selected sections on propositional and first-order logic.;%
Genesereth and Nilsson, \textit{Logical Foundations of Artificial Intelligence}, selected chapters.%
}

\newcommand{\AIMLUnitIIISources}{%
Mueller and Massaron, \textit{Machine Learning for Dummies}, For Dummies, 2016.;%
Mitchell, \textit{Machine Learning}, selected chapters on learning basics and evaluation.;%
Scikit-learn documentation, accessed 2026-02-28.%
}

\newcommand{\AIMLUnitIVSources}{%
Mueller and Massaron, \textit{Machine Learning for Dummies}, For Dummies, 2016.;%
James, Witten, Hastie, and Tibshirani, \textit{An Introduction to Statistical Learning} (2e), selected chapters.;%
Scikit-learn documentation, accessed 2026-02-28.%
}

\newcommand{\AIMLUnitVSources}{%
NIST, \textit{AI Risk Management Framework (AI RMF 1.0)}, 2023.;%
Barocas, Hardt, and Narayanan, \textit{Fairness and Machine Learning} (online book), selected sections.;%
Knaflic, \textit{Storytelling with Data}, selected chapters on communicating results.%
}
%
}{%
  % Faculty copies live under `private/` so their compile CWD is different.
  \IfFileExists{../../../../public/_shared/aiml-sources.tex}{%
    % Source strings for Elements of AIML (used by slides + notes).

\newcommand{\AIMLRepoURL}{https://github.com/tali7c/Elements-of-AIML}

\newcommand{\AIMLLabSources}{%
Provost and Fawcett, \textit{Data Science for Business}.;%
McKinney, \textit{Python for Data Analysis}.;%
WEKA Documentation (Waikato Environment for Knowledge Analysis), accessed 2026-02-28.;%
Matplotlib and Seaborn documentation, accessed 2026-02-28.%
}

\newcommand{\AIMLUnitISources}{%
Rich and Knight, \textit{Artificial Intelligence}, McGraw Hill, 2017.;%
Russell and Norvig, \textit{Artificial Intelligence: A Modern Approach} (4e), Ch. 1--2 (Introduction, Intelligent Agents).;%
Poole and Mackworth, \textit{Artificial Intelligence: Foundations of Computational Agents} (2e), selected sections.%
}

\newcommand{\AIMLUnitIISources}{%
Russell and Norvig, \textit{Artificial Intelligence: A Modern Approach} (4e), Ch. 7--9 (Logical Agents, First-Order Logic, Inference).;%
Huth and Ryan, \textit{Logic in Computer Science} (2e), selected sections on propositional and first-order logic.;%
Genesereth and Nilsson, \textit{Logical Foundations of Artificial Intelligence}, selected chapters.%
}

\newcommand{\AIMLUnitIIISources}{%
Mueller and Massaron, \textit{Machine Learning for Dummies}, For Dummies, 2016.;%
Mitchell, \textit{Machine Learning}, selected chapters on learning basics and evaluation.;%
Scikit-learn documentation, accessed 2026-02-28.%
}

\newcommand{\AIMLUnitIVSources}{%
Mueller and Massaron, \textit{Machine Learning for Dummies}, For Dummies, 2016.;%
James, Witten, Hastie, and Tibshirani, \textit{An Introduction to Statistical Learning} (2e), selected chapters.;%
Scikit-learn documentation, accessed 2026-02-28.%
}

\newcommand{\AIMLUnitVSources}{%
NIST, \textit{AI Risk Management Framework (AI RMF 1.0)}, 2023.;%
Barocas, Hardt, and Narayanan, \textit{Fairness and Machine Learning} (online book), selected sections.;%
Knaflic, \textit{Storytelling with Data}, selected chapters on communicating results.%
}
%
  }{%
    % Source strings for Elements of AIML (used by slides + notes).

\newcommand{\AIMLRepoURL}{https://github.com/tali7c/Elements-of-AIML}

\newcommand{\AIMLLabSources}{%
Provost and Fawcett, \textit{Data Science for Business}.;%
McKinney, \textit{Python for Data Analysis}.;%
WEKA Documentation (Waikato Environment for Knowledge Analysis), accessed 2026-02-28.;%
Matplotlib and Seaborn documentation, accessed 2026-02-28.%
}

\newcommand{\AIMLUnitISources}{%
Rich and Knight, \textit{Artificial Intelligence}, McGraw Hill, 2017.;%
Russell and Norvig, \textit{Artificial Intelligence: A Modern Approach} (4e), Ch. 1--2 (Introduction, Intelligent Agents).;%
Poole and Mackworth, \textit{Artificial Intelligence: Foundations of Computational Agents} (2e), selected sections.%
}

\newcommand{\AIMLUnitIISources}{%
Russell and Norvig, \textit{Artificial Intelligence: A Modern Approach} (4e), Ch. 7--9 (Logical Agents, First-Order Logic, Inference).;%
Huth and Ryan, \textit{Logic in Computer Science} (2e), selected sections on propositional and first-order logic.;%
Genesereth and Nilsson, \textit{Logical Foundations of Artificial Intelligence}, selected chapters.%
}

\newcommand{\AIMLUnitIIISources}{%
Mueller and Massaron, \textit{Machine Learning for Dummies}, For Dummies, 2016.;%
Mitchell, \textit{Machine Learning}, selected chapters on learning basics and evaluation.;%
Scikit-learn documentation, accessed 2026-02-28.%
}

\newcommand{\AIMLUnitIVSources}{%
Mueller and Massaron, \textit{Machine Learning for Dummies}, For Dummies, 2016.;%
James, Witten, Hastie, and Tibshirani, \textit{An Introduction to Statistical Learning} (2e), selected chapters.;%
Scikit-learn documentation, accessed 2026-02-28.%
}

\newcommand{\AIMLUnitVSources}{%
NIST, \textit{AI Risk Management Framework (AI RMF 1.0)}, 2023.;%
Barocas, Hardt, and Narayanan, \textit{Fairness and Machine Learning} (online book), selected sections.;%
Knaflic, \textit{Storytelling with Data}, selected chapters on communicating results.%
}
%
  }%
}


\title{Elements of AIML\\Miscellaneous -- Lecture 02 Notes\\Eigenvalues and Eigenvectors from a Data Sample}
\author{Tofik Ali}
\date{\today}

\begin{document}
\maketitle
\tableofcontents

\section{Lecture Overview}
This supplementary lecture shows the full manual path from a tiny \textbf{raw data sample} to:
\begin{itemize}[leftmargin=*]
  \item the mean vector,
  \item the centered data matrix,
  \item the covariance matrix,
  \item the eigenvalues,
  \item the eigenvectors,
  \item the PCA interpretation of those eigenvectors.
\end{itemize}
This is useful because students often understand the definitions but do not see how the calculation is chained from the original sample.
\notesources{\AIMLMiscSources}

\section{How to use these notes}
Use this lecture after:
\begin{itemize}[leftmargin=*]
  \item \textbf{Unit 03 Lecture 03}, where PCA is introduced conceptually,
  \item \textbf{Miscellaneous Lecture 01}, if you want to know why PCA and LDA are tied to eigen-analysis.
\end{itemize}
The present lecture is not mainly about comparing methods; it is about doing one derivation carefully and correctly.

\section{Learning Outcomes}
\begin{itemize}[leftmargin=*]
  \item Compute the mean vector from a small dataset.
  \item Center the data correctly.
  \item Build the covariance matrix from the centered data.
  \item Solve the characteristic equation to obtain eigenvalues.
  \item Solve the linear system for eigenvectors.
  \item Interpret the largest-eigenvalue direction as the first principal component.
\end{itemize}

\section{Raw sample data}
Consider the following three data points:
\[
  x_1 =
  \begin{bmatrix}
    2\\1
  \end{bmatrix},
  \qquad
  x_2 =
  \begin{bmatrix}
    1\\2
  \end{bmatrix},
  \qquad
  x_3 =
  \begin{bmatrix}
    0\\0
  \end{bmatrix}
\]
These are deliberately small so the entire calculation can be done by hand.

\section{Step 1: Compute the mean vector}
The sample mean is:
\[
  \mu =
  \frac{1}{3}(x_1 + x_2 + x_3)
\]
Substitute the values:
\[
  \mu =
  \frac{1}{3}
  \left(
  \begin{bmatrix}
    2\\1
  \end{bmatrix}
  +
  \begin{bmatrix}
    1\\2
  \end{bmatrix}
  +
  \begin{bmatrix}
    0\\0
  \end{bmatrix}
  \right)
  =
  \frac{1}{3}
  \begin{bmatrix}
    3\\3
  \end{bmatrix}
  =
  \begin{bmatrix}
    1\\1
  \end{bmatrix}
\]

\section{Step 2: Center each sample}
Centering means subtracting the mean vector from every sample:
\[
  x_i^{(c)} = x_i - \mu
\]
So:
\[
  x_1^{(c)} =
  \begin{bmatrix}
    2\\1
  \end{bmatrix}
  -
  \begin{bmatrix}
    1\\1
  \end{bmatrix}
  =
  \begin{bmatrix}
    1\\0
  \end{bmatrix}
\]
\[
  x_2^{(c)} =
  \begin{bmatrix}
    1\\2
  \end{bmatrix}
  -
  \begin{bmatrix}
    1\\1
  \end{bmatrix}
  =
  \begin{bmatrix}
    0\\1
  \end{bmatrix}
\]
\[
  x_3^{(c)} =
  \begin{bmatrix}
    0\\0
  \end{bmatrix}
  -
  \begin{bmatrix}
    1\\1
  \end{bmatrix}
  =
  \begin{bmatrix}
    -1\\-1
  \end{bmatrix}
\]

Therefore the centered data matrix is:
\[
  X_c =
  \begin{bmatrix}
    1 & 0\\
    0 & 1\\
    -1 & -1
  \end{bmatrix}
\]
where each row is one centered sample.

\section{Step 3: Compute the covariance matrix}
For $n=3$ samples, the sample covariance matrix is:
\[
  C = \frac{1}{n-1} X_c^\top X_c = \frac{1}{2} X_c^\top X_c
\]

\subsection{Compute $X_c^\top X_c$}
\[
  X_c^\top =
  \begin{bmatrix}
    1 & 0 & -1\\
    0 & 1 & -1
  \end{bmatrix}
\]
Now multiply:
\[
  X_c^\top X_c =
  \begin{bmatrix}
    1 & 0 & -1\\
    0 & 1 & -1
  \end{bmatrix}
  \begin{bmatrix}
    1 & 0\\
    0 & 1\\
    -1 & -1
  \end{bmatrix}
\]
\[
  =
  \begin{bmatrix}
    1^2 + 0^2 + (-1)^2 & 1\cdot 0 + 0\cdot 1 + (-1)(-1)\\
    0\cdot 1 + 1\cdot 0 + (-1)(-1) & 0^2 + 1^2 + (-1)^2
  \end{bmatrix}
  =
  \begin{bmatrix}
    2 & 1\\
    1 & 2
  \end{bmatrix}
\]
So:
\[
  C =
  \frac{1}{2}
  \begin{bmatrix}
    2 & 1\\
    1 & 2
  \end{bmatrix}
  =
  \begin{bmatrix}
    1 & 0.5\\
    0.5 & 1
  \end{bmatrix}
\]

\section{Step 4: Find the eigenvalues}
We solve:
\[
  \det(C - \lambda I)=0
\]
Substitute $C$:
\[
  \det
  \begin{bmatrix}
    1-\lambda & 0.5\\
    0.5 & 1-\lambda
  \end{bmatrix}
  = 0
\]
So:
\[
  (1-\lambda)^2 - (0.5)(0.5) = 0
\]
\[
  (1-\lambda)^2 - 0.25 = 0
\]
\[
  1 - 2\lambda + \lambda^2 - 0.25 = 0
\]
\[
  \lambda^2 - 2\lambda + 0.75 = 0
\]
Factor:
\[
  (\lambda - 1.5)(\lambda - 0.5)=0
\]
Hence the eigenvalues are:
\[
  \lambda_1 = 1.5,
  \qquad
  \lambda_2 = 0.5
\]

\section{Step 5: Find the eigenvectors}
\subsection{Eigenvector for $\lambda_1 = 1.5$}
Compute:
\[
  C - 1.5I =
  \begin{bmatrix}
    1 & 0.5\\
    0.5 & 1
  \end{bmatrix}
  -
  \begin{bmatrix}
    1.5 & 0\\
    0 & 1.5
  \end{bmatrix}
  =
  \begin{bmatrix}
    -0.5 & 0.5\\
    0.5 & -0.5
  \end{bmatrix}
\]
Now solve:
\[
  \begin{bmatrix}
    -0.5 & 0.5\\
    0.5 & -0.5
  \end{bmatrix}
  \begin{bmatrix}
    x\\y
  \end{bmatrix}
  = 0
\]
This gives:
\[
  -0.5x + 0.5y = 0
  \Rightarrow y=x
\]
So one eigenvector is:
\[
  v_1 =
  \begin{bmatrix}
    1\\1
  \end{bmatrix}
\]

\subsection{Eigenvector for $\lambda_2 = 0.5$}
Compute:
\[
  C - 0.5I =
  \begin{bmatrix}
    0.5 & 0.5\\
    0.5 & 0.5
  \end{bmatrix}
\]
Now solve:
\[
  \begin{bmatrix}
    0.5 & 0.5\\
    0.5 & 0.5
  \end{bmatrix}
  \begin{bmatrix}
    x\\y
  \end{bmatrix}
  = 0
\]
This gives:
\[
  0.5x + 0.5y = 0
  \Rightarrow y=-x
\]
So one eigenvector is:
\[
  v_2 =
  \begin{bmatrix}
    1\\-1
  \end{bmatrix}
\]

\section{Step 6: Normalize the eigenvectors}
To convert them into unit vectors:
\[
  \|v_1\| = \sqrt{1^2 + 1^2} = \sqrt{2},
  \qquad
  \|v_2\| = \sqrt{1^2 + (-1)^2} = \sqrt{2}
\]
Therefore:
\[
  \hat v_1 =
  \frac{1}{\sqrt{2}}
  \begin{bmatrix}
    1\\1
  \end{bmatrix},
  \qquad
  \hat v_2 =
  \frac{1}{\sqrt{2}}
  \begin{bmatrix}
    1\\-1
  \end{bmatrix}
\]

\section{Interpretation}
Because:
\[
  \lambda_1 = 1.5 > \lambda_2 = 0.5
\]
the direction:
\[
  \hat v_1 =
  \frac{1}{\sqrt{2}}
  \begin{bmatrix}
    1\\1
  \end{bmatrix}
\]
captures more variance than the perpendicular direction.

\subsection{PCA meaning}
If we use this covariance matrix in PCA:
\begin{itemize}[leftmargin=*]
  \item $\hat v_1$ becomes the first principal component,
  \item $\hat v_2$ becomes the second principal component,
  \item the first component is more informative because its eigenvalue is larger.
\end{itemize}

\section{What this example teaches}
This worked example shows a full chain:
\begin{enumerate}[leftmargin=*]
  \item start from raw samples,
  \item compute the mean,
  \item center the samples,
  \item build the covariance matrix,
  \item solve the characteristic equation,
  \item solve for the eigenvectors,
  \item interpret the result in PCA language.
\end{enumerate}

\section{Common mistakes}
\begin{itemize}[leftmargin=*]
  \item \textbf{Skipping centering:} then the covariance matrix does not reflect spread around the mean correctly.
  \item \textbf{Using the raw data matrix directly:} in PCA the eigen-problem is applied to the covariance matrix, not to the raw sample matrix.
  \item \textbf{Ignoring normalization:} direction is enough mathematically, but unit vectors make the interpretation cleaner.
  \item \textbf{Confusing largest eigenvalue with largest coordinate:} the eigenvalue measures variance along the direction, not the magnitude of one coordinate.
\end{itemize}

\section{Practice Questions (with answers)}
\subsection*{Q1. Why do we divide by $n-1$ in the covariance formula?}
\textbf{Answer:} because we are computing the sample covariance, which uses Bessel's correction.

\subsection*{Q2. Why is the covariance matrix symmetric?}
\textbf{Answer:} because covariance between feature 1 and feature 2 equals covariance between feature 2 and feature 1, so $C = C^\top$.

\subsection*{Q3. Which eigenvector becomes PC1?}
\textbf{Answer:} the eigenvector corresponding to the largest eigenvalue.

\subsection*{Q4. Why is this example useful for PCA but not enough for ICA?}
\textbf{Answer:} because PCA depends directly on covariance eigen-analysis, while ICA requires a stronger independence-based optimization beyond covariance alone.

\section{Summary}
For a tiny dataset, eigenvalues and eigenvectors can be found by hand.
This derivation is important because it makes the PCA procedure less mysterious:
\begin{itemize}[leftmargin=*]
  \item covariance summarizes feature spread,
  \item eigenvectors give the principal directions,
  \item eigenvalues tell how much variance those directions capture.
\end{itemize}
\notesources{\AIMLMiscSources}

\end{document}
